% Complete whole-H rational-coordinate proof and arithmetic application.
% Independent review: greenberg_magnus_adaptation.md and sibling audit.
% Include cup_ext_pole_section.tex before this section.
\subsection{Greenberg's one-variable hypothesis suffices}
\label{sec:greenberg-kernel}

The integral freeness and cyclicity used in the preceding pole argument
can be replaced by their versions after inverting $p$. The change must
be made in the coefficient algebra of the entire free pro-$p$ group:
a subgroup with finite index on abelianizations need not be open.

\begin{theorem}[K17-GREENBERG-KERNEL]\label{thm:greenberg-kernel}
Let $p$ be an odd prime, put $F_r=\mathbb Q(\mu_{p^r})$ and
$F_\infty=\bigcup_r F_r$, and let
\[
 X=\varprojlim_r\operatorname{Cl}(F_r)\{p\}
\]
with transition maps given by norms. If $X^+$ is finite, then
\[
 \Omega_p^*=\Omega_p.
\]
The hypothesis is the one-variable CM form of Greenberg's conjecture.
No simplicity of Iwasawa roots is assumed, and no converse is asserted.
\end{theorem}

We first prove the rational-coordinate version of pole separation and
the precise cohomological comparison needed for its arithmetic use.

\begin{lemma}[K17-RATIONAL-MAGNUS]\label{lem:rational-magnus}
Let $p$ be odd and let $H$ be a second-countable free pro-$p$ group
with commuting continuous actions of
$\Gamma=\langle\gamma\rangle\simeq\mathbb Z_p$ and
$\Delta\simeq\mathbb F_p^\times$. Let $c\in\Delta$ have order two,
and put
\[
 \Lambda=\mathbb Z_p[[\Gamma]],\qquad R=\Lambda[1/p],\qquad
 M=H^{\rm ab}=U\oplus T,
\]
where $c=-1$ on $U$ and $c=1$ on $T$. Assume that $M$ is finitely
generated over $\Lambda$, that $U[1/p]$ is free of rank one over $R$
in every odd $\Delta$-character, and that $T$ is finite free over
$\mathbb Z_p$, with every nonzero $T_e[1/p]$ cyclic over $R$.
Assume also that abelian eigenclasses have exact $\Delta$-eigenlifts,
including inside arbitrarily small $\Delta$-stable neighbourhoods.

Let $\rho:H\to\mathcal U(\mathbb Q_p)$ be continuous and
$\Gamma\times\Delta$-equivariant, where $\mathcal U$ is the
weight-completed free prounipotent group on letters $s_m$ in all odd
weights $m\ge3$. Suppose its image is Zariski dense in every finite
weight quotient and, after a Levi normalization,
\[
 \gamma s_m=u^m s_m,\qquad \delta s_m=\omega(\delta)^m s_m,
\]
where $u$ topologically generates $1+p\mathbb Z_p$.

For each nonzero $T_e$, choose an integral multiple $v_e$ of a rational
cyclic generator, an exact eigenlift $y_e\in H$, and its actual
conjugates $y_{e,j}=\gamma^j y_e\gamma^{-j}$. Let $q_e(z)$ be the
characteristic polynomial of $\gamma$ on $T_e$. In a rational
$R$-basis of $U[1/p]$, let $C_{e;a,b}(T_1,T_2)$ be the odd--odd
quadratic coefficients of the actual companion error. Suppose that
for every root $\alpha$ of every $q_e$, there is a matching ordered
odd pair $(a,b)$ for which
\[
 C_{e;a,b}\left(T,\frac{\alpha}{1+T}-1\right)\ne0.
                                                               \tag{GM.1}
\]
Then $\rho$ is injective. All coefficient series here have a common
finite denominator in each fixed assigned degree; no bound uniform
over all degrees is required.
\end{lemma}

\begin{proof}
\emph{An actual compact basis.}
A compact torsionfree $\mathbb Z_p$-module is a product of copies of
$\mathbb Z_p$: its Pontryagin dual is a divisible $p$-primary group,
hence a direct sum of copies of $\mathbb Q_p/\mathbb Z_p$.
The abelianization of free $H$ is torsionfree and metrizable.
Using the integral character idempotents, choose pointed convergent
$\mathbb Z_p$-bases in the $U_a$, and finite bases in the $T_e$.
The open quotient map $H\to M$ permits lifts of the convergent
basis vectors tending to $1$. Apply the exact-eigenlift construction
inside a nested sequence of $\Delta$-stable open neighbourhoods;
it preserves both the abelian classes and their convergence.

Let $F$ be the free pro-$p$ group on this pointed sequence and the
finite even set. Its abelianization is the product module on those
basis vectors, so the resulting map $F\to H$ is an isomorphism on
abelianizations. Its compact image is Frattini-surjective and hence
equals $H$. Indeed a proper closed subgroup of any pro-$p$ group is
contained in a proper open subgroup and then in an index-$p$ subgroup,
so cannot be Frattini-surjective. Since $H$ is free, the surjection
has a continuous section. The section is inverse on abelianizations;
its closed image is Frattini-surjective in $F$ and is all of $F$.
Thus $F\simeq H$. This argument does not use an infinite-rank
Hopficity assertion.

Give the odd Magnus variables $X_i$ assigned degree $1$, and the
finite even variables $Y_l$ assigned degree $2$. This Magnus
expansion is injective. To see this directly, kill a tail of the
convergent free sequence. Every finite quotient kills such a tail,
so $H$ is the inverse limit of these finite-rank free groups.
An ordinary finite-rank Magnus expansion detects every nonidentity
element in some finite assigned degree. At each fixed degree the
coefficient module is a finite sum of completed tensor products of
$U$ and $T$, using these assigned degrees.

\emph{Odd rational coordinates.}
Choose a character-preserving isomorphism
\[
 U[1/p]\simeq\bigoplus_{a\ {
m odd}}Re_a.
\]
Finite generation over $\Lambda$ clears the denominators of this
isomorphism and its inverse. Thus there are continuous integral
maps between $U$ and $\Lambda^{(p-1)/2}$ whose two compositions are
multiplication by the same power of $p$. Their completed tensor
powers have compositions equal to the corresponding powers of that
scalar. The tensor modules are torsionfree: in the product bases
they are products of copies of $\mathbb Z_p$. Hence the maps are
injective and become inverse isomorphisms after inverting $p$.
Every fixed degree therefore has coordinates in
\[
 \mathbb Z_p[[\Gamma^s]][1/p]
   =\mathbb Z_p[[T_1,\ldots,T_s]][1/p],                       \tag{GM.2}
\]
for its singleton positions and their character types. These are
the Amice transforms of bounded measures on $\Gamma^s$.
The denominator bound may depend on the degree, which causes no
problem at a first nonzero degree.

The target linear coefficient of weight $m$ is a continuous additive
map from $M$ to $\mathbb Q_p$. Parity kills $T$, and character
equivariance kills every $U_a$ except $a\equiv m\pmod{p-1}$.
On the remaining rank-one $R$-module, $\Gamma$-equivariance makes
it a scalar multiple of evaluation at $T=u^m-1$: check this on
polynomials in $\gamma$ and then pass to bounded series by
continuity. The scalar is nonzero by Zariski density, since the
weight-$m$ abelian coordinate of the target cannot vanish on a
dense image. Rescale $s_m$ by this scalar. Such weightwise rescaling
is an automorphism of the completed free Lie algebra, because each
finite weight quotient has finitely many generators. The singleton
kernel is now exactly evaluation at $u^m-1$ in its matching type.

\emph{Even companion coordinates.}
Put $\lambda_e=\operatorname{rank}_{\mathbb Z_p}T_e$. The classes
\[
 v_e,\gamma v_e,\ldots,\gamma^{\lambda_e-1}v_e
\]
are a rational basis, though they need not be an integral basis.
The polynomial $q_e$ is monic integral, has unit constant term,
and reduces to $(z-1)^{\lambda_e}$ modulo $p$, since the image of
$\Gamma$ in $\operatorname{GL}(T_e/p)$ is a finite $p$-group.
It is also the minimal polynomial, by rational cyclicity.

The actual conjugates $y_{e,j}$ have assigned degree at least $2$
and target generator length at least $2$. Their assigned-degree-two
parts have the form
\[
 Z_{e,j}=\sum_l b_{j,l}Y_{e,l}+A_{e,j}(X),                    \tag{GM.3}
\]
where $b$ is a finite invertible rational matrix and $A_{e,j}$ is
a completed quadratic odd tensor. Before the odd coordinate change
these coefficients are integral; afterwards they retain bounded
denominators. There are no other terms of assigned degree two.
The quadratic tensor is alternating, by the commutator description
or primitiveness of the first Magnus coefficient.

Replace the finite variables $Y$ by $Z$. This is a homogeneous
triangular automorphism of the completed coefficient algebra:
it fixes $X$ and its inverse is $Y=b^{-1}(Z-A(X))$.
In each fixed degree only finitely many even variables occur,
so both directions preserve a common finite denominator bound.
In particular a first nonzero degree remains nonzero.
This is a change in the rational tensor algebra, not a claim that
the companion classes form an integral group basis.

By definition of $Z_{e,j}$, its leading target is the length-two
part of the actual element $y_{e,j}$. For an ordered target pair
$(m,n)$, put $z=u^{m+n}$. Its coefficient satisfies
\[
 B_{e,j;m,n}=z^jB_{e,0;m,n}.
\]
Write $q_e(z)=z^{\lambda_e}-\sum_jc_{e,j}z^j$ and form
\[
 E_e=(\gamma^{\lambda_e}y_e\gamma^{-\lambda_e})
       \left(\prod_jy_{e,j}^{c_{e,j}}\right)^{-1}.
\]
Its abelian class is zero, so $E_e$ is in the closed derived
subgroup. Its assigned-degree-two part is the odd alternating
tensor $C_e$. There are no target length-one terms in any even
factor, so comparison of length two gives exactly
\[
 q_e(z)B_{e,0;m,n}
   =C_{e;a,b}(u^m-1,u^n-1),\qquad
 a\equiv m,\ b\equiv n\pmod{p-1}.                           \tag{GM.4}
\]
This is the same companion class as the actual extension in
Lemma~\ref{lem:cup-pole-bridge}. Exact eigenlifts make its assigned
degree-two character $e$; hence wrong-character pairs have zero
coefficient. One can also obtain this zero from (GM.4) on the full
residue grid and the successive analytic identity theorem.

\emph{Partial fractions and continuity.}
Pass to a finite splitting field $L$, with integers $\mathcal O$.
All roots $\alpha$ satisfy $\alpha-1\in\mathfrak m_{\mathcal O}$.
The finite-dimensional rational-function space with basis
$z^j/q_e(z)$ has alternative basis
\[
 (z-\alpha)^{-k},\qquad
 1\le k\le\operatorname{mult}_e(\alpha).
\]
The corresponding constant invertible change in the finite $Z$
variables preserves assigned degree and bounded denominators in
every fixed degree. Its even letter $(e,\alpha,k)$ has kernel
\[
 \frac{C_{e;a,b}(T_1,T_2)}
      {((1+T_1)(1+T_2)-\alpha)^k}.                          \tag{GM.5}
\]
The numerator is common to all letters from the same companion
component, even when roots repeat.

There are finitely many positive integers $w$ such that $u^w$ is a
root of some $q_e$. Restrict all test weights to a sufficiently
large ordinary tail, excluding every vanishing pair denominator.
In each fixed odd residue class this still includes sequences
$m_0+(p-1)p^j$ with an interior $p$-adic accumulation point, for
a sufficiently large fixed representative $m_0$.

Substitution is legitimate at every finite arithmetic-weight cutoff.
In the actual compact basis, truncate its convergent sequence and
apply finite-rank Magnus substitution. These truncations converge
to the identity on $H$, since every finite group quotient kills a
tail. The image of $H$ in the finite target quotient is compact,
so its coefficient functions are bounded and continuous. Passing
to the limit therefore gives the completed tensor substitution.
The bounded rational coordinate changes preserve this statement in
each fixed degree. In (GM.2) it is equivalently integration against
bounded measures and the finite Fubini rule. We do not invoke a
free pro-$p$ universal property for a noncompact rational group.

\emph{Separation of a first nonzero degree.}
Suppose $g\ne1$ and $\rho(g)=1$, and let $d$ be its first nonzero
assigned degree after these invertible changes. A larger source
degree cannot contribute to target length $d$. Its coefficient is
a finite sum over ordered tilings of $d$ positions by singleton
odd types and pairs labelled $(e,\alpha,k)$:
\[
 \begin{gathered}
 \sum_{\mathcal P}A_{\mathcal P}(T_i:\ i\text{ singleton})
       \prod_{(j,e,\alpha,k)\in\mathcal P}
       \frac{C_{e;a_j,a_{j+1}}(T_j,T_{j+1})}
            {D_{j,\alpha}^{k}}=0,\\
 D_{j,\alpha}=(1+T_j)(1+T_{j+1})-\alpha.
 \end{gathered}                                            \tag{GM.6}
\]
Every $A_{\mathcal P}$ has a common bounded denominator. Clear
the finitely many denominators. Vanishing on the Cartesian grid
of large odd test weights and the successive one-variable identity
theorem make (GM.6) an identity in the fraction field of
$\mathcal O[[T_1,\ldots,T_d]][1/p]$.

Select a nonzero pattern with the largest number of pairs. Fix its
pair positions $J$ and its roots and even characters, and among
nonzero patterns with those data choose a componentwise maximal
pole-order vector $(k_j)$. At each selected pair choose the ordered
odd types in (GM.1), and at each singleton its source type. These
choices are independent because the pairs are disjoint.

Set $S_j=D_{j,\alpha_j}$ for $j\in J$. Retain as base variables
the singleton variables and the left variable in each selected pair.
The substitutions
\[
 T_{j+1}=\frac{\alpha_j+S_j}{1+T_j}-1
\]
are automorphisms of the complete integral power-series ring.
Let $K_0$ be the fraction field of the bounded-denominator base
ring. The identity embeds in $K_0[[S_j:j\in J]][S_j^{-1}:j\in J]$,
with finite lower bounds on exponents. All other denominators are
units in $K_0[[S_j]]$: a different root at the same position has
nonzero constant root difference; a different pair position joins
two distinct selected blocks, whose surviving variables are
independent, so its product is not identically a prescribed root.
Wrong even characters have already vanished.

Take the coefficient of $\prod_{j\in J}S_j^{-k_j}$.
A pattern missing a selected position or using a different root
there has no negative exponent in that variable. A pattern
containing all selected pairs and an extra pair has zero coefficient
by the first maximality choice. Among remaining patterns, a
contribution requires all pole orders to dominate $k_j$; the
componentwise maximal choice excludes strict domination and an
incomparable vector fails in at least one variable. Singleton types
must match. Thus exactly one coefficient remains:
\[
 A_{\mathcal P}(T_i:\ i\text{ singleton})
 \prod_{j\in J}
 C_{e_j;a_j,b_j}\left(T_j,\frac{\alpha_j}{1+T_j}-1\right).
\]
The first factor is nonzero and its coefficient ring injects into
$K_0$; every other factor is nonzero by (GM.1). Their product is
nonzero in a field, a contradiction. If $J$ is empty, the same
argument is the ordinary singleton moment identity. This proves
injectivity of $\rho$.
\end{proof}

\begin{lemma}[K17-ETALE-PROP]\label{lem:etale-prop}
Let $S$ be a connected normal affine scheme, let
$G=\pi_1^{\mathrm{\acute et}}(S)$, and let $P$ be its maximal
pro-$p$ quotient. For a finite $p$-primary local coefficient module
$A$ whose action factors through $P$, the natural maps give
\[
 H^1(P,A)\simeq H^1_{\mathrm{\acute et}}(S,A),\qquad
 H^2(P,A)\hookrightarrow H^2_{\mathrm{\acute et}}(S,A).
\]
They are compatible with cup products and with finite external
automorphisms. Consequently generation of a character summand of
the target $H^2$ by cups of $H^1$ classes gives the same generation
in group cohomology and makes the displayed injection an
isomorphism on that summand.
\end{lemma}

\begin{proof}
Put $N=\ker(G\to P)$. An element of $H^1(N,A)^P$ is a continuous
$P$-equivariant homomorphism $N\to A$, since $N$ acts trivially
on $A$. If its image were nonzero, its kernel would be $G$-normal
and the quotient of $G$ by that kernel would be an extension of
$P$ by a nonzero finite $p$-group. This is a pro-$p$ quotient:
every finite quotient has $p$-power order. It would be strictly
larger than the maximal pro-$p$ quotient, a contradiction.
Thus $H^1(N,A)^P=0$. The five-term Hochschild--Serre sequence gives
\[
 H^1(P,A)\simeq H^1(G,A),\qquad
 H^2(P,A)\hookrightarrow H^2(G,A).
\]
No finite-generation hypothesis on $N$ was used.

The full universal finite \'etale tower over $S$ has zero $H^1$
with finite locally constant coefficients: after trivializing the
coefficients, an $H^1$ torsor is finite \'etale and becomes trivial
at a later stage of the same tower. Finite \'etale objects and
torsors descend under these affine limits. The Cartan--Leray
spectral sequence for the tower therefore gives
\[
 H^1(G,A)\simeq H^1_{\mathrm{\acute et}}(S,A),\qquad
 H^2(G,A)\hookrightarrow H^2_{\mathrm{\acute et}}(S,A).
\]
This is only its low-degree consequence, and asserts no vanishing
of the higher rows. The sequence, the profinite-tower passage and
the simply connected $H^1$ vanishing are given explicitly in
\cite[Theorem~14.9, Remark~14.10, Example~14.11 and
Proposition~14.13, pp.~96--98]{MilneLEC2013}.
All comparison maps are natural and respect cups.
Lift each cup generator using the $H^1$ isomorphism and add its
group cups; this proves surjectivity onto the asserted summand.
Injectivity then identifies every group class in that summand with
the lifted sum.
\end{proof}

\begin{proof}[Proof of Theorem~\ref{thm:greenberg-kernel}]
The prime $3$ is already covered by the regular-prime result, so
assume $p\ge5$. Put
\[
 P=\operatorname{Gal}(\Omega_p/F_1),\quad
 H=\operatorname{Gal}(\Omega_p/F_\infty),\quad
 \Gamma=\operatorname{Gal}(F_\infty/F_1),\quad
 M=H^{\rm ab}=U\oplus T.
\]
Choose the finite $\Delta$ lift by the coprime complement theorem.
The fibre over a generator of $\Gamma$ is a $\Delta$-stable
$H$-torsor because the cyclotomic quotient is abelian.
The coprime identity $H^1(\Delta,H)=1$, obtained from the finite
complement theorem and compactness, gives a $\Delta$-fixed lift
$\gamma$. Its procyclic closure maps isomorphically to $\Gamma$
and supplies the commuting section used above.
The arithmetic structure used earlier, from Iwasawa and
Ferrero--Washington, gives free pro-$p$ $H$.
The finite-level Kummer sequence makes $H^1(P,\mathbb F_p)$
finite dimensional, so $P$ is finitely generated and $H$ is
second-countable. Adjust a finite generating set of $P$ by powers
of the chosen $\Gamma$ section to obtain elements of $H$.
Together with $\gamma$ they generate $P$; their closed normal
closure is $H$, since the remaining quotient is the given
procyclic quotient. Their abelian classes therefore generate $M$
over $\Lambda=\mathbb Z_p[[\Gamma]]$.
The actual-image results of Hain--Matsumoto and Brown give the
continuous dense rational representation and Levi normalization in
Lemma~\ref{lem:rational-magnus}; this is the same actual
representation as in the preceding pole theorem, not an assumed
faithful motivic representation. The exact eigenlift construction
is \cite[Lemma~2.1, pp.~3--4]{Sharifi2002}; performing it inside
stable neighbourhoods preserves convergence. We verify the changed
module and cup hypotheses in full.

\emph{Rational odd module.}
The ring $R=\Lambda[1/p]$ is a PID. Indeed $\Lambda$ is a
two-dimensional regular local UFD, and its only height-two prime
contains $p$. Thus $R$ is a one-dimensional noetherian UFD and
hence a PID. It remains the ring of bounded-denominator series,
not $\mathbb Q_p[[T]]$.

For an odd character $a\ne\omega,\omega^{-1}$, apply the full
duality calculation in Wake--Wang-Erickson's Proposition~2.2.6
with $\theta=\omega a^{-1}$. It gives
\[
 \operatorname{Ext}^1_\Lambda(M_a^{\#}(1),\Lambda)=X_\theta,
 \qquad
 \operatorname{Ext}^2_\Lambda(M_a^{\#}(1),\Lambda)=0.           \tag{GM.7}
\]
Here $M$ is their $p$-ramified module $\mathfrak X$, whereas $X$
on the right is unramified. Since $\theta$ is even, Greenberg
makes the right-hand side zero after inverting $p$. Ext commutes
with this localization for finitely generated modules. Over a PID,
the Ext group into the ring of a finitely generated torsion summand
is nonzero unless that summand is zero. Thus $M_a[1/p]$ is free.
Its rank is one by the local/global unit exact sequence in the
same proof: local units have rank one in this character, global
minus units contribute only roots of unity, and the unramified
quotient is torsion. The inversion $\#$ and Tate twist transport
the ring by automorphisms, so undoing them preserves freeness
\cite[Proposition~2.2.6 and proof, pp.~3--4]{WakeWangErickson2017}.

The two excluded characters can be calculated directly.
Coleman's exact sequence at tame level one is
\[
 0\longrightarrow\mathbb Z_p(1)\longrightarrow U_{\rm loc}
 \longrightarrow\mathbb Z_p[[\Delta\times\Gamma]]
 \longrightarrow\mathbb Z_p(1)\longrightarrow0
\]
\cite[Theorem~5.5.31, p.~148; Proposition~6.1.1,
pp.~151--152]{SharifiIwasawaNotes}.
Global minus units are exactly the root-of-unity term, and
Stickelberger gives $X_\omega=X_{\omega^{-1}}=0$; the latter
exceptional statement is explicitly recorded in
\cite[proof of Lemma~5.2.7, p.~13]{WakeWangErickson2017}.
The $\omega^{-1}$ component of $M$ is consequently free of rank
one. In the $\omega$ component, quotienting the local units by
their root-of-unity kernel identifies $M_\omega$ with
\[
 \ker(\Lambda\longrightarrow\mathbb Z_p(1))
     =(\gamma-u)\Lambda,
\]
again free of rank one. The two characters are distinct for
$p\ge5$. The valuation direction at the unique prime above $p$
is even and does not affect this calculation. We have therefore
proved the required rational freeness for every odd character.

\emph{Rational even module.}
Kummer duality identifies $T=M^+$ up to finite kernel and cokernel
with $(X^-)^\iota(1)$
\cite[Proposition~3.4.7 and Corollary~3.4.8,
p.~92]{SharifiIwasawaNotes}.
Ferrero--Washington makes it finitely generated over $\mathbb Z_p$,
and freeness of $H$ excludes $\mathbb Z_p$-torsion. Thus $T$ is
finite free. Greenberg implies that $X^-$ is pseudocyclic
character by character, by the exact local/global unit sequence and
Iwasawa adjunction in
\cite[Corollary~2.1.3 and Conjecture~2.1.5 with the following
explanation, p.~8]{WakeWangEricksonPseudo2018}.
Finite kernels and cokernels vanish after inverting $p$, so every
$T_e[1/p]$ is cyclic. This allows repeated roots; it does not assert
integral cyclicity without Vandiver.

\emph{Integral symbol surjectivity under Greenberg.}
Fix an even character $\theta\ne1,\omega^2$ and put
$\chi=\omega^{-1}\theta$. Wake--Wang-Erickson's
Corollary~8.1.3(2) assumes exactly
\[
 \begin{gathered}
 X_\chi(1)[1/p]\text{ cyclic over }(\Lambda/\xi_\chi)[1/p],\\
 X_\theta\otimes_\Lambda(\Lambda/\xi_\chi)[1/p]=0.
 \end{gathered}                                            \tag{GM.8}
\]
The localization in the second condition is essential.
Both conditions follow from the just stated consequences of
Greenberg. The conclusion is an integral isomorphism
\[
 \varpi:(\mathcal H^-_\theta/I_\theta\mathcal H^-_\theta)
           /(p\text{-primary torsion})
       \xrightarrow{\ \sim\ }X_\chi(1),                    \tag{GM.9}
\]
so the original $\varpi$ is surjective
\cite[Corollary~8.1.3 and complete proof,
pp.~48--49]{WakeWangEricksonPseudo2018}.
The source proves this with repeated roots: the derivative identity
makes the error nilpotent in each cyclic local factor, rather than
requiring the derivative to be invertible. Thus (GM.9) does not
insert a simple-root hypothesis.

Fukaya--Kato, \S7.1.10, deduce from this surjectivity that for every
$r\ge1$ the finite-level $\theta$ component is generated by the
symbols
\[
 \{1-\zeta_{p^r}^{a},1-\zeta_{p^r}^{b}\}.
\]
Their \S7.1.9 identifies its reduction modulo $p^r$ with the
corresponding class-group quotient, and their \S6.1.5 shows that
at tame level one the only excluded even characters are
$1,\omega^2$
\cite[\S6.1.5, p.~76; \S\S7.1.8--7.1.10,
p.~86]{FukayaKato2012}.
Those excluded characters correspond to the unramified components
$X_{\omega^{-1}}$ and $X_\omega$, which are zero. Norm projections
from $X^-$ onto each finite minus class group are surjective, so
the omitted finite targets are zero as well
\cite[Proposition~3.4.1 and proof,
pp.~90--91]{SharifiIwasawaNotes}.
Components with unit characteristic power series are also zero:
the source's $X^-$ is $p$-torsionfree, and a finite such module must
vanish. We apply the nonunit running assumption of the modular
source only to the nonzero components. Hence every even target is
covered.

\emph{All finite group levels and coefficient conventions.}
Let
\[
 P_n=H\rtimes\Gamma^{p^n},\qquad
 S_n=\operatorname{Spec}\mathcal O_{F_{n+1}}[1/p].
\]
The maximal pro-$p$ quotient of $\pi_1(S_n)$ is exactly $P_n$.
Indeed the normal closure over $F_1$ of any finite Galois
$p$-extension of $F_{n+1}$ unramified outside $p$ remains a
$p$-extension: its conjugate compositum over $F_{n+1}$ has
$p$-group Galois group, and $F_{n+1}/F_1$ has $p$-power degree.
Thus it lies in the fixed field $\Omega_p$ used throughout.

At $F_{n+1}$, the finite-level symbol theorem gives generation
with coefficients $p^{n+1}$. Reduction gives generation with
coefficients $p^n$, exactly as in Lemma~\ref{lem:cup-pole-bridge}.
The class-group quotient description in FK \S7.1.9 makes this
reduction explicit. The Brauer contribution is zero: these fields
have no real places, and only one finite prime is omitted, so the
global invariant condition forces its lone possible local
invariant to vanish.

For a $p$-unit $v$, the quotient $v/c(v)$ is a global unit of
absolute value one at every embedding and hence is a root of
unity. Consequently the negative $c$ part of the $p$-unit group
modulo $p^n$ is cyclic. Its alternating self-cup is zero because
$2$ is a unit; mixed-parity cups lie in the odd target. Thus the
even target is generated by cups of $c$-even unit classes. This
argument does not identify the whole of $H^1$ with units and does
not use Vandiver.

Apply Lemma~\ref{lem:etale-prop}. Each Kummer class has a unique
lift to $H^1(P_n,\mathbb Z/p^n(1))$, and group $H^2$ injects into
\'etale $H^2$. Lifting the even symbol generators makes this
injection surjective on the even part and gives precisely
\[
 H^2(P_n,\mathbb Z/p^n(2))^+
  =\bigl\langle f\smile g\mid
       f,g\in H^1(P_n,\mathbb Z/p^n(1))^+\bigr\rangle.       \tag{GM.10}
\]
This treats every $n\ge1$; no compatibility of expressions as sums
of cups is needed. Scalar extension to the integers of a finite
splitting field preserves the assertion, since that ring is finite
free over $\mathbb Z_p$.

Lemma~\ref{lem:cup-pole-bridge} now gives the nonzero companion
restriction (GM.1) for every root of every even component, including
repeated roots. Its dual root is $u^2/\alpha$ and its dual character
is $\omega^{2-e}$, so the inversion and both Tate twists match
the symbol target. Lemma~\ref{lem:rational-magnus} proves
injectivity on the entire $H$.

Finally an element in the kernel of the original pro-$p$ tripod
action is cyclotomically trivial, hence lies in $H$. Functorial
rational completion sends it to $1$ under $\rho$, and the proved
injectivity makes it trivial. This is exactly
$\Omega_p^*=\Omega_p$.
\end{proof}

The theorem is conditional only on the stated one-variable
Greenberg assertion. In particular it does not assert Greenberg
for unverified primes. Vandiver supplies the stronger condition
$X^+=0$ by the earlier persistence proof, so
Theorem~\ref{thm:vandiver-kernel} and its verified numerical range
remain unconditional consequences of their stated computations.
